\section{Phát kiến nằm ở mục tiêu, không nằm ở kiến trúc}
\label{sec:ch09-phat-kien}

\index{tiền huấn luyện!hai pha}%
Mọi mô hình trong sách này cho tới đây đều huấn luyện một lần, trên dữ liệu của
đúng bài toán nó phải giải, với nhãn của đúng bài toán ấy. Corpus của
chương~\ref{ch:encoder-decoder} là các cặp câu Anh-Việt, và mô hình học dịch.

Hai bài của mục này bẻ đôi chuyện đó. \tn{Tiền huấn luyện}{pretraining} là pha
một: huấn luyện trên văn bản trơn không có nhãn nào. \tn{Tinh
chỉnh}{fine-tuning} là pha hai: lấy mô hình ấy và chỉnh nó sang một bài toán có
nhãn, với ít dữ liệu hơn nhiều bậc.

GPT-1 viết cả hai pha ra thành công thức. Phương trình (1) là mục tiêu pha một,
đoán token tiếp theo. Phương trình (4) là mục tiêu pha hai, đoán nhãn. Phương
trình (5) ghép chúng lại, \(L_3(C) = L_2(C) + \lambda L_1(C)\), và mục 4.1 đặt
\(\lambda\) bằng 0.5, tức là pha một không tắt hẳn trong pha hai mà ở lại làm
một số hạng \tn{điều chuẩn}{regularization}.

Cái mua được bằng phép bẻ đôi ấy là dữ liệu. Nhãn thì đắt và hiếm; văn bản trơn
thì không. BooksCorpus mà GPT-1 dùng có \enquote{over 7,000 unique unpublished
books}, và bài nói rõ vì sao chọn nó thay vì một corpus cùng cỡ:
\enquote{Crucially, it contains long stretches of contiguous text, which allows
the generative model to learn to condition on long-range
information}~\autocite{radford2018gpt1}. BERT dùng BooksCorpus 800 triệu từ cộng
Wikipedia tiếng Anh 2500 triệu từ~\autocite{devlin2019bert}. GPT-2 dựng WebText,
\enquote{slightly over 8 million documents for a total of 40 GB of
text}~\autocite{radford2019gpt2}.

Không tập nào trong ba tập ấy có nhãn. Đó là toàn bộ phát kiến, và nó là một
phát kiến về mục tiêu huấn luyện chứ không phải về kiến trúc.

\subsection*{Bỏ mask thì mất mục tiêu, nên phải tìm mục tiêu khác}

\index{tiền huấn luyện!masked language model}%
Mục~\ref{sec:ch09-cung-mot-khoi} đã nói chỗ này: bỏ mask nhân quả đi thì mục
tiêu \enquote{đoán token tiếp theo} sụp, vì token tiếp theo nằm trong tầm nhìn
của chính vị trí phải đoán nó. Mục~\ref{sec:ch07-che} đo được điều đó trên
corpus của sách bằng phép can thiệp, và bài BERT phát biểu nó bằng lời ở mục
3.1.

Cái BERT đặt vào chỗ trống là masked language model. Thủ tục, nguyên văn:

\begin{quote}
\enquote{In all of our experiments, we mask 15\% of all WordPiece tokens in each
sequence at random. [...] If the i-th token is chosen, we replace the i-th token
with (1) the [MASK] token 80\% of the time (2) a random token 10\% of the time
(3) the unchanged i-th token 10\% of the time.}
\end{quote}

Ba nhánh 80-10-10 tồn tại vì token \code{[MASK]} chỉ xuất hiện lúc tiền huấn
luyện và không bao giờ xuất hiện lúc tinh chỉnh. Nếu mọi vị trí phải đoán đều
mang \code{[MASK]}, mô hình học được rằng chỉ cần chú ý tới các vị trí có dấu
ấy, và pha hai đưa cho nó một câu không có dấu nào. Hai nhánh 10\% còn lại là để
mô hình không biết chắc vị trí nào đang bị hỏi.

\subsection*{Cái giá của 15 phần trăm, tính bằng số học}

\index{tiền huấn luyện!mật độ tín hiệu của hai mục tiêu}%
Con số 15 ấy có một cái giá mà tôi chưa thấy bài nào trong hai bài nêu ra, và
nó tính được mà không cần chạy gì.

Một lượt truyền xuôi qua một chuỗi \(n\) token tốn chừng ấy phép tính dù mục
tiêu là gì: kiến trúc không biết nó đang được huấn luyện kiểu nào. Nhưng số
\emph{đích} mà lượt ấy sinh ra thì khác hẳn. Một \tn{mô hình ngôn ngữ}{language
model} trái sang phải chấm điểm mọi vị trí, vì mỗi vị trí là một bài đoán hợp lệ
với prefix của riêng nó. Masked language model chỉ chấm các vị trí bị che.

\begin{measured}{mật độ tín hiệu của hai mục tiêu, trên cùng một lượt truyền}
\begin{minted}{text}
objective                targets/token  FLOPs/target, relative
left-to-right LM         1.0000         1.0000
masked LM at 15%         0.1500         6.6667

at BERT-base's shape, one pretraining sequence of 512 tokens:
  positions in the sequence        512
  positions an LM predicts         512
  positions BERT predicts          76
\end{minted}

Số học thuần túy từ 15\% của mục 3.1 và độ dài 512 của phụ lục A.2 trong bài
BERT. Đo bằng \code{experiments/ch09\_counts.py} tại tag \code{ch09}.
\end{measured}

Trên cùng một chuỗi và cùng một lượt truyền, BERT thu về 76 đích còn một mô
hình ngôn ngữ thu về 512. Tính theo phép tính bỏ ra cho mỗi đích, mục tiêu che
đắt hơn 6.67 lần.

Đó không phải một lời chê bài BERT. Đổi lại, mỗi đích của BERT được nhìn cả hai
phía, còn mỗi đích của mô hình ngôn ngữ chỉ nhìn được bên trái, và bài đo được
điều đó có giá. Bảng 5 mục 5.1 chạy \tn{phép cắt bỏ}{ablation}:

\begin{quote}
\enquote{We first examine the impact brought by the NSP task. In Table 5, we
show that removing NSP hurts performance significantly on QNLI, MNLI, and SQuAD
1.1.}
\end{quote}

Hàng \code{LTR \& No NSP} của bảng ấy là hàng đáng đọc: bỏ cả hai chiều lẫn
next-sentence prediction thì Stanford Question Answering Dataset, viết tắt là
SQuAD, tụt từ 88.5 xuống 77.8 F1, trong khi bỏ
riêng next-sentence prediction chỉ làm nó xuống 87.9. Phần lớn khoảng cách là
của tính hai chiều chứ không của mục tiêu câu.

Nhưng chỗ đánh đổi giữa \enquote{nhiều đích một chiều} và \enquote{ít đích hai
chiều} thì không bài nào trong hai bài đo, vì hai bài không so với nhau trên
cùng một ngân sách tính toán. Ngành sau này quay lại đúng câu hỏi ấy, và
bài tập cuối chương có một câu hỏi ở đó.

\subsection*{Hai pha, và chỗ chi phí dồn vào}

\index{tiền huấn luyện!bất đối xứng chi phí hai pha}%
Con số làm cho cả câu chuyện này có nghĩa nằm ở phụ lục A.2 và mục 4 của bài
BERT, và hai câu ấy nên đọc liền nhau.

Pha một: \enquote{Training of BERT\_BASE was performed on 4 Cloud TPUs in Pod
configuration (16 TPU chips total). Training of BERT\_LARGE was performed on 16
Cloud TPUs (64 TPU chips total). Each pre-training took 4 days to complete.}

Pha hai: \enquote{Compared to pre-training, fine-tuning is relatively
inexpensive. All of the results in the paper can be replicated in at most 1 hour
on a single Cloud TPU, or a few hours on a GPU, starting from the exact same
pre-trained model.}

Bốn ngày trên 64 chip, một lần. Một giờ trên một chip, cho mỗi bài toán, và bài
đo trên một loạt bài toán. Cái bất đối xứng ấy là lý do hai pha đáng bẻ đôi: chi
phí lớn trả một lần rồi chia cho mọi bài toán về sau, và mỗi bài toán mới chỉ
phải trả phần nhỏ.

Tóm tắt của bài nói cái đó bằng một câu về kiến trúc, và bây giờ thì câu ấy đọc
ra được: \enquote{the pre-trained BERT model can be fine-tuned with just one
additional output layer [...] without substantial task-specific architecture
modifications.} Một tầng ra thêm vào. Không phải một kiến trúc mới cho mỗi bài
toán, mà là một kiến trúc cộng một tầng.

GPT-1 giải cùng bài toán ấy theo cách khác và nói ra ở mục 3.3: thay vì đổi mô
hình cho vừa bài toán, nó đổi bài toán cho vừa mô hình, bằng cách viết đầu vào
có cấu trúc thành một chuỗi tuyến tính. \enquote{These input transformations
allow us to avoid making extensive changes to the architecture across tasks.}

Hai bài, hai cách, cùng một mục đích: giữ cho phần đắt tiền dùng lại được.
